\documentclass[10pt,a4paper]{article}
\usepackage{amsmath,amssymb,amsthm,graphicx}
\usepackage{amsmath,amssymb,amsthm,hyperref,xfrac}
\providecommand{\ie}{i.e.\ }
\providecommand{\skpt}{\smallskip}
\usepackage[scr=rsfs,cal=euler]{mathalfa}
\multlinegap0pt

\hyphenation{begin decom-po-si-tion sequence}
\newenvironment{enumeratea}
{\bgroup\def\theenumi{\alph{enumi}}\begin{enumerate}}
{\end{enumerate}\egroup}

\newcommand{\ba}{{\boldsymbol a}}
\newcommand{\bb}{{\boldsymbol b}}
\newcommand{\bd}{{\boldsymbol d}}
\newcommand{\bu}{{\boldsymbol u}}
\newcommand{\bx}{{\boldsymbol x}}
\newcommand{\bmy}{{\boldsymbol y}}

\newcommand{\bD}{{\boldsymbol D}}
\newcommand{\bL}{{\boldsymbol L}}
\newcommand{\bP}{{\boldsymbol P}}
\newcommand{\bV}{{\boldsymbol V}}
\newcommand{\bX}{{\boldsymbol X}}
\newcommand{\bZ}{{\boldsymbol Z}}

\newcommand{\Psfrac}[2]{(\sfrac{#1}{#2})}
\newcommand{\psfrac}[2]{\sfrac{(#1)}{#2}}
\newcommand{\bigO}[1]{O(#1)}

\usepackage{mathtools}
\usepackage{bbm}

\newcommand{\NN}{\mathbb{N}}
\newcommand{\ZZ}{\mathbb{Z}}
\newcommand{\RR}{\mathbb{R}}
\newcommand{\CC}{\mathbb{C}}
\newcommand{\PP}{\mathbb{P}}
\newcommand{\EE}{\mathbb{E}}

\newcommand{\boldalpha}{\boldsymbol{\alpha}}
\newcommand{\bolddelta}{\boldsymbol{\delta}}
\newcommand{\boldrho}{\boldsymbol{\rho}}
\newcommand{\dee}{\,\mathrm{d}}
\newcommand{\one}{\mathbbm{1}}
\newcommand{\BC}{\mathcal{B}}
\newcommand{\EC}{\mathcal{E}}
\newcommand{\RS}{\mathscr{R}}
\newcommand{\prob}[1]{\PP[ #1 ]}
\newcommand{\bigoo}[1]{O\!\left(#1 \right)}
\newcommand{\abs}[1]{\left\lvert #1 \right\rvert}
\newcommand{\floor}[1]{\left\lfloor #1 \right\rfloor}
\newcommand{\riemint}[4]{\int_{#1}^{#2}\! #3 \, \mathrm{d}#4}
\newcommand{\maxnorm}[1]{\left\| #1 \right\|_{\infty}}

\renewcommand{\Re}{\textup{Re}}
\renewcommand{\Im}{\textup{Im}}

\newcommand{\dtv}{\mathrm{d}_{\mathrm{TV}}}
\newcommand{\dw}{\mathrm{d}_{\mathrm W}}

\let\epsilon\varepsilon

\newcommand{\pextra}{P_{\mathtt{extra}}}
\newcommand{\parratia}{P_{\mathtt{Arratia}}}

\newcommand{\simplex}{\Delta^{k-1}}

\newcommand{\bg}{\big}
\newcommand{\bgg}{\Big}
\newcommand{\bggg}{\bigg}

\newcommand{\Rnt}{R_{\mathrm{NT}}}
\newcommand{\Rprob}{R_{\mathrm{PD}}}
\newcommand{\ECnt}{\EC_{\mathrm{NT}}}
\newcommand{\ECprob}{\EC_{\mathrm{PD}}}

\newcommand{\Bnt}{\mathcal{B}_{\mathrm{NT}}}
\newcommand{\Bprob}{\mathcal{B}_{\mathrm{PD}}}

\DeclareMathOperator{\vard}{Var}
\DeclareMathOperator{\Dir}{Dir}

\theoremstyle{plain}
\newtheorem{cor}{Corollary}[section]
\newtheorem{lem}[cor]{Lemma}
\newtheorem{prop}[cor]{Proposition}
\newtheorem{conj}[cor]{Conjecture}
\newtheorem{theo}{Theorem}

\theoremstyle{definition}
\newtheorem{example}{Example}
\newtheorem{deff}[cor]{Definition}
\newtheorem{rem}{Remark}[section]
\newtheorem*{rems}{Remark}
\newtheorem*{remss}{Remarks}

\makeatletter
\ams@newcommand{\multiint}[1]{\DOTSI\protect\MultiIntegral{#1}}
\renewcommand{\MultiIntegral}[1]{%
\edef\ints@c{\noexpand\intop
\ifnum#1=\z@\noexpand\intdots@\else\noexpand\intkern@\fi
\replicate{#1-2}{\noexpand\intop\noexpand\intkern@}%
\noexpand\intop
\noexpand\ilimits@
}%
\futurelet\@let@token\ints@a
}
\makeatother
\ExplSyntaxOn
\cs_new:Npn \replicate #1 #2 { \prg_replicate:nn { #1 } { #2 } }
\ExplSyntaxOff
\title{A maximal coupling of two laws on the positive integers: Y = f(X;U,U') with P[X != Y] = d\_TV(mu,nu) and Y of law nu}
\author{Tony Haddad, Dimitris Koukoulopoulos}
\date{doi:10.5802/jep.317; CC BY 4.0}
\begin{document}
\maketitle
\noindent\textit{Editorial note.} The lemma is attributed to Arratia in the source heading; the authors reproduce his proof at 1767-1793 and that reproduction is the proof carried here.\par\medskip
\noindent\textit{Editorial note.} The supporting counting-measure identity 1715-1737 is included because the main proof invokes it at 1783; it is labelled as a supporting obligation and no separate count is assigned to it.\par\medskip
\noindent\textit{Editorial note.} The map f\_\{mu,nu\} is the same function the paper later uses to extract N\_x (line 488); the context anchor 440-441 records that role without importing any part of the coupling theorem.\par\medskip
\noindent\textit{Editorial note.} All displayed statements are taken from the fixed publisher source unchanged.\par\medskip
\noindent\textit{Allowed standard prerequisites (editorial).}\begin{itemize}
\item Total variation distance of two laws on the positive integers. d\_TV(mu,nu) = sup\_A |mu(A) - nu(A)|, and for laws supported on N it equals sum\_i (mu(i)-nu(i))\textasciicircum{}+ = sum\_i (mu(i)-nu(i))\textasciicircum{}-.. Supplies the quantity the constructed coupling attains exactly; the identity form is proved inside this card rather than assumed.
\item Uniform randomisation on the unit square. Given X of law mu and two independent uniforms U, U' on (0,1), the two-branch map f\_\{mu,nu\}(X;U,U') changes value exactly when U*mu(X) > nu(X).. Converts the total variation identity into an explicit coupling with Y measurable with respect to (X,U,U'); verifying the law of Y is the second obligation.
\end{itemize}\medskip
\section{The coupling}
\label{sec:coupling}

The \textit{total variation distance} is a metric between two probability distributions. Let $\mu$ and $\nu$ be two probability measures on $\CC$. Then
\begin{equation}
\label{eq:dtv dfn}
\dtv(\mu, \nu) \coloneqq \sup |\mu(A) - \nu(A)|,
\end{equation}
where the supremum is taken over all Lebesgue-measurable subsets of $\CC$. For any real number $a$, let
\[
a^+ \coloneqq \max\{a, 0\} \quad \text{and} \quad a^- \coloneqq \max\{-a, 0\}.
\]When $\mu$ and $\nu$ are supported on $\ZZ_{\ge 1}$, here are some alternative definitions of the total variation distance:

The total variation distance will be especially useful in this paper because of the following proposition. In~\cite[\S 3.8]{Arratia02}, Arratia proved that for any two random variables~$X$ and~$Y$ returning positive integers, we~can always construct $X'$ and $Y'$ within the same probability space such that:
\begin{itemize}
\item $X'$ and $Y'$ have the same distribution as $X$ and $Y$, respectively;

\item $Y'$ is a function of $X', U, V$ where $(U, V)$ is a point uniformly chosen in the unit square independent of $X'$;

\item $\PP[X' \ne Y'] = \dtv(X, Y)$.
\end{itemize}

\begin{lem}[Arratia, \cite{Arratia02}]
\label{lem:coupling-dtv}
Let $\mu$ and $\nu$ be two probability measures supported on~$\NN$, let $X$ be a random variable with law $\mu$, and let $U$ and $U'$ be two uniform random variables in $(0, 1)$ such that $X, U, U'$ are independent. Let $Y \coloneqq f_{\mu, \nu}(X; U, U')$ with $f_{\mu, \nu}$ defined as above. Then $\PP[X \ne Y] = \dtv(\mu, \nu)$ and $\PP[Y \in A] = \nu(A)$.
\end{lem}

\begin{lem}
\label{lem:equiv-def-dtv}
Let $\mu$ and $\nu$ be two probability measure supported on $\NN$. Then
\[
\dtv(\mu, \nu) = \sum_{i \ge 1} (\mu(i) - \nu(i))^+ = \sum_{i \ge 1} (\mu(i) - \nu(i))^-.
\]
\end{lem}

\begin{proof}
Let $E \coloneqq \{i \in \NN \, : \, \mu(i) > \nu(i)\}$ and let $\xi\coloneqq \mu - \nu$. Note that for any $B\subseteq \NN$, we~have $\xi(B\cap E) \ge 0 \ge \xi(B\cap E^c)$. Therefore, for any $A \subseteq \NN$, we~have
\[
\xi(A) = \xi(E) + \xi(A \cap E^c) - \xi(A^c \cap E) \le \xi(E)
\]
and
\[
\xi(A) = \xi(E^c) + \xi(A \cap E) - \xi(A^c \cap E^c) \ge \xi(E^c).
\]
Therefore, $\dtv(\mu, \nu) = \sup_{A \subseteq \NN} |\xi(A)| = \max\{\xi(E), -\xi(E^c)\}$. Note that $\xi(E) = \sum_{i \ge 1} (\mu(i) - \nu(i))^+$, that $-\xi(E^c) = \sum_{i \ge 1} (\mu(i) - \nu(i))^-$ and that
\[
\sum_{i \ge 1} (\mu(i) - \nu(i))^+ = \sum_{i \ge 1} \bg[(\mu(i) - \nu(i)) + (\mu(i) - \nu(i))^-\bg] = \sum_{i \ge 1} (\mu(i) - \nu(i))^-.
\]
The lemma follows.
\end{proof}

We repeat his proof here. Let $\mu$ and $\nu$ be two probability measures supported on $\NN$. Let $z_j \coloneqq \sum_{i \le j} \frac{(\mu(i) - \nu(i))^-}{\dtv(\mu, \nu)}$ (with $z_0 \coloneqq 0$). We~consider the function $f_{\mu, \nu}\colon \NN \times (0, 1)^2 \to \NN$ defined as
\[
f_{\mu, \nu}(m; a, b) \coloneqq \begin{cases}
m &\text{if $a \cdot \mu(m) \le \nu(m)$,} \\
\sum_{i \ge 1} i \cdot \one_{z_{i-1} < b \le z_i} &\text{otherwise.}
\end{cases}
\]
This is the function used in Section \ref{sec:coupling} for the extraction of $N_x$.
Note that if $a,b\in(0,1)$ and $m\in\NN$, then we have the equivalency
\begin{equation}
\label{eq:f TV property}
f_{\mu, \nu}(m; a, b) \ne m\quad\iff\quad a\cdot \mu(m) > \nu(m).
\end{equation}
Indeed, the direction ``$\Rightarrow$'' is obvious. To see the converse direction, note that if \hbox{$a\cdot\mu(m)>\nu(m)$}, then $(\mu(m)-\nu(m))^-=0$, and hence the interval $(z_{m-1},z_m]$ is empty.

\begin{proof}
Using \eqref{eq:f TV property}, we~find that
\begin{equation}
\label{eq:event X neq Y}
\{X \ne Y\}=\{U \cdot \mu(X) > \nu(X)\}.
\end{equation}
Furthermore, we~directly compute that
\begin{equation}
\label{eq:prob X eq m and U large}
\PP\bg[U\cdot \mu(m)> \nu(m), X = m\bg] = (\mu(m) - \nu(m))^+.
\end{equation}
Therefore,
\begin{equation}
\label{eq:prob X neq Y}
\PP[X \ne Y] = \PP[U \cdot \mu(X) > \nu(X)] = \sum_{m \ge 1} (\mu(m) - \nu(m))^+ = \dtv(\mu, \nu)
\end{equation}
with Lemma \ref{lem:equiv-def-dtv}.

Next, we~prove that $\PP[Y = n] = \nu(n)$ for any $n \in \NN$. Note that $Y = n$ if, and only if, one of two disjoint events happen: either we have $U \cdot \mu(n)\le \nu(n)$ with $X = n$, or we have $U\cdot \mu(X) > \nu(X)$ with $z_{n-1} < U' \le z_n$. Therefore, with \eqref{eq:prob X eq m and U large} and \eqref{eq:prob X neq Y}, we~have
\begin{align*}
\PP[Y = n]
&= \PP\bg[U \cdot \mu(n)\le \nu(n), X = n\bg] + \PP\bg[U\cdot \mu(X) > \nu(X)\bg] \cdot \PP\bg[z_{n-1} < U' \le z_n\bg] \\
&= \bgg(\mu(n) - \bg(\mu(n) - \nu(n)\bg)^+\bgg) + \dtv(\mu, \nu) \cdot \frac{(\mu(n) - \nu(n))^-}{\dtv(\mu, \nu)} = \nu(n),
\end{align*}
where we used \eqref{eq:event X neq Y} and \eqref{eq:prob X neq Y} to show that $\PP\bg[U\cdot \mu(X) > \nu(X)\bg]=\dtv(\mu,\nu)$.
This concludes the proof of the lemma.
\end{proof}

\begin{thebibliography}{99}
\bibitem{Arratia98} R. Arratia - “On the central role of scale invariant Poisson processes on (0,∞)”, in Microsurveys in discrete probability (Princeton, NJ, 1997), DIMACS Ser. Discrete Math. Theoret. Comput. Sci., vol. 41, American Mathematical Society, Providence, RI, 1998, p. 21-41
\bibitem{Arratia02} R. Arratia - “On the amount of dependence in the prime factorization of a uniform random integer”, in Contemporary combinatorics, Bolyai Soc. Math. Stud., vol. 10, János Bolyai Math. Soc., Budapest, 2002, p. 29-91
\bibitem{ABT06} R. Arratia, A. D. Barbour \& S. Tavaré - “A tale of three couplings: Poisson-Dirichlet and GEM approximations for random permutations”, Combin. Probab. Comput. 15 (2006) no. 1-2, p. 31-62
\bibitem{BareikisManstavicius07} G. Bareikis \& E. Manstavičius - “On the DDT theorem”, Acta Arith. 126 (2007) no. 2, p. 155-168
\bibitem{Bill72} P. Billingsley - “On the distribution of large prime divisors”, Period. Math. Hungar. 2 (1972), p. 283-289
\bibitem{delaBretecheTenenbaum16} R. de la Bretèche \& G. Tenenbaum - “Sur les processus arithmétiques liés aux diviseurs”, Adv. in Appl. Probab. 48 (2016) no. A, p. 63-76
\bibitem{DDT79} J.-M. Deshouillers, F. Dress \& G. Tenenbaum - “Lois de répartition des diviseurs. I”, Acta Arith. 34 (1979) no. 4, p. 273-285 (loose errata)
\bibitem{DonnellyTavare87} P. Donnelly \& S. Tavaré - “The population genealogy of the infinitely-many neutral alleles model”, J. Math. Biol. 25 (1987) no. 4, p. 381-391
\bibitem{Feng10} S. Feng - The Poisson-Dirichlet distribution and related topics. Models and asymptotic behaviors, Probability and its Appl., Springer, Heidelberg, 2010
\bibitem{JKB} N. L. Johnson, S. Kotz \& N. Balakrishnan - Discrete multivariate distributions, Wiley Series in Probab. and Stat.: Applied Probab. and Stat., John Wiley \& Sons, Inc., New York, 1997
\bibitem{Kingman93} J. F. C. Kingman - Poisson processes, Oxford Studies in Probab., vol. 3, The Clarendon Press, Oxford University Press, New York, 1993
\bibitem{KotzBal00} S. Kotz, N. Balakrishnan \& N. L. Johnson - Continuous multivariate distributions. Vol. 1. Models and applications, Wiley Series in Probab. and Stat.: Applied Probab. and Stat., Wiley-Interscience, New York, 2000
\bibitem{Koukoulo19} D. Koukoulopoulos - The distribution of prime numbers, Graduate Studies in Math., vol. 203, American Mathematical Society, Providence, RI, 2019
\bibitem{LastPenrose18} G. Last \& M. Penrose - Lectures on the Poisson process, Institute of Math. Stat. Textbooks, vol. 7, Cambridge University Press, Cambridge, 2018
\bibitem{Ken22} S.-K. Leung - “Dirichlet law for factorisation of integers, polynomials and permutations”, Math. Proc. Cambridge Philos. Soc. 175 (2023) no. 3, p. 649-676
\bibitem{KenWeb24} S.-K. Leung - “Dirichlet law with arbitrary parameters for factorizations of integers” (2025), Retrieved on 24 August 2025 from https://sites.google.com/view/sunkaileung/miscellaneous
\bibitem{MV07} H. L. Montgomery \& R. C. Vaughan - Multiplicative number theory. I. Classical theory, Cambridge Studies in Advanced Math., vol. 97, Cambridge University Press, Cambridge, 2007
\bibitem{NyandwiSmati13} S. Nyandwi \& A. Smati - “Distribution laws of pairs of divisors”, Integers 13 (2013), article ID A13, 13 pages
\bibitem{Tenenbaum00} G. Tenenbaum - “A rate estimate in Billingsley’s theorem for the size distribution of large prime factors”, Q. J. Math. 51 (2000) no. 3, p. 385-403
\end{thebibliography}
\end{document}
